%!TEX root = ../../csuthesis_main.tex
\chapter{绪论}
\section{研究背景与意义}
\subsection{研究背景}
随着自动驾驶仿真、驾驶模拟器、数字孪生技术快速发展，车辆动力学仿真成为虚拟测试的核心基础。传统游戏引擎内置车辆仿真大多采用纯刚体射线悬架模型，轮胎、悬架全部视为刚体，忽略轮胎形变、柔性部件变形，仿真真实度有限，难以满足工程级仿真需求。
UE4 采用 PhysX3 作为底层物理，UE5 替换为 Chaos 物理求解器，二者原生车辆模块以刚体动力学为主，缺少对软体、柔性部件的仿真支持。NVIDIA PhysX5 作为开源多物理仿真 SDK，新增有限元 FEM 软体动力学、基于位置动力学 PBD两套可 GPU 加速的可变形体求解框架，除传统刚体车辆动力学之外，还支持柔性体、布料、颗粒、可充气物体仿真，能够对轮胎形变、车身柔性件实现物理建模，为高保真车辆仿真提供了新的技术路径。
现有研究中，PhysX 更多用于刚体破碎、工业设备演示，如圆锥破碎机设备仿真、预破碎岩石仿真；而车辆仿真大多只使用 PhysX 传统 Vehicle 刚体模块，很少将 FEM/PBD 软体框架引入车辆轮胎、柔性悬架部件仿真当中。
本课题基于 PhysX5 开源 SDK，结合 FEM 有限元软体动力学与 PBD 基于位置动力学，构建一套高保真车辆仿真系统，实现车辆直行、转弯、碰撞等交互式仿真场景，探究刚度、阻尼等物理参数对仿真真实感、运行实时性的影响。课题研究成果可以为驾驶模拟器、自动驾驶虚拟测试、车辆数字孪生提供技术参考。
\subsubsection{研究意义}
理论意义：探究 PhysX5 中 FEM、PBD 两套软体框架在车辆部件仿真中的适配性，对比不同刚度、阻尼参数下仿真真实度与计算开销，补充 PhysX 软体动力学在车辆仿真领域的应用研究。
实践意义：搭建可交互车辆仿真原型系统，不仅实现传统刚体车辆动力学，同时引入柔性轮胎、柔性部件仿真，完成直行、转向、碰撞测试场景，可为虚拟驾驶、教学演示、车辆预研仿真提供原型参考。

\section{国内外研究现状}
\subsection{国内研究现状}
国内车辆仿真领域，大量研究基于Carla、Prescan、UE4/UE5等平台开展，底层物理求解大多依托PhysX3、Chaos刚体动力学，重点面向自动驾驶感知、路径规划、车辆控制算法验证，仿真对象关注车辆运动轨迹，忽略车体与轮胎的柔性形变。
在软体动力学方向，国内不少学者将PBD算法进行二次开发移植到虚幻引擎，多用于软组织手术仿真、柔性机器人、布料形变仿真，验证了PBD求解器实时仿真的优势。但现有研究大多针对独立柔性部件，很少将PhysX5原生FEM、PBD模块用于整车一体化仿真，缺少针对车辆场景的参数对照实验，缺少对刚度、阻尼参数如何平衡保真度与实时性的系统性分析。总体来看，国内对于PhysX5软体动力学在高保真车辆仿真中的应用仍处于探索阶段。

\subsection{国外研究现状}
国外NVIDIA团队持续维护PhysX开源项目，在Omniverse平台中不断完善PBD、FEM求解器，公开了软体碰撞、约束求解、并行加速的底层技术文档，为柔性物体实时仿真提供了底层支撑。国际上已有不少学术工作将PBD、FEM用于轮胎、柔性杆件等单个车辆构件仿真，验证了软体模型相比刚体模型在形变还原上的优势。
但现有国外相关工作依然存在局限：多数研究只针对轮胎、防撞梁等单一柔性零部件，很少搭建包含整车模型、转向驱动逻辑、碰撞交互的完整仿真场景；同时相关文献较少设置多组梯度参数对照，缺少定量对比分析刚度、阻尼对仿真真实感和计算开销的影响，整车级柔性交互仿真案例仍然偏少。

\subsection{研究述评}
综合国内外现有研究成果可以发现：
第一，现阶段主流车辆仿真仍然以刚体动力学为主，仿真重心在运动控制，难以满足高保真形变仿真需求；
第二，PhysX5已经具备成熟的PBD、FEM软体求解能力，在柔性物体仿真上已经得到验证，但在整车仿真场景下的落地应用并不充分；
第三，现有研究缺少面向车辆场景的参数对比实验，对于刚度、阻尼参数如何取舍、两种软体求解器各自的优劣，缺少量化结论；
第四，缺少一套轻量化、可交互，集成直行、转弯、碰撞场景的PhysX5柔性车辆仿真Demo。

\section{研究内容与创新点}
本文以PhysX5物理引擎为基础，利用引擎原生的FEM有限元软体动力学、PBD基于位置动力学框架，完成高保真车辆仿真系统开发。主要研究内容如下：
1. 梳理PhysX5刚体、PBD、FEM软体动力学基础理论，对比PhysX3、Chaos与PhysX5软体模块的差异；
2. 完成仿真系统总体方案设计，建立车辆软体模型，设计驱动、转向控制逻辑，搭建直行、转弯、障碍物碰撞三类交互场景；
3. 设计多组对照实验，改变刚度、阻尼等核心物理参数，从视觉真实感、运行帧率两个维度记录结果，对比PBD和FEM方案的表现差异；
4. 完成仿真程序调试、结果整理与分析，总结参数选取的经验规律。
本文创新点如下：
1. 区别于UE4 PhysX3、UE5 Chaos以刚体约束为主的车辆仿真方案，直接使用PhysX5原生PBD/FEM软体动力学实现整车柔性仿真，还原车体、轮胎的形变效果；
2. 面向车辆仿真场景设计梯度参数对照实验，定量分析刚度、阻尼参数对仿真保真度与实时性能的影响，对比两种软体求解框架的适用场景；
3. 搭建集成直行、转向、碰撞的一体化交互仿真场景，实现轻量化可复现的车辆柔性仿真Demo。

\section{论文组织架构}
本文一共分为五个章节，各章节内容安排如下：
第一章为绪论。阐述课题的研究背景与研究意义，梳理国内外车辆仿真、PhysX软体动力学相关研究现状并进行研究述评，明确本文的研究内容与创新点，介绍全文组织结构。
第二章为相关理论与技术基础。介绍PhysX5引擎整体架构，对比PhysX3、Chaos的差异；阐述PBD基于位置动力学、FEM有限元软体动力学的基本原理，说明软体仿真核心参数刚度、阻尼的物理含义。
第三章为仿真系统方案设计。明确系统总体架构，完成车辆模型划分、软体网格构建、驱动与转向控制逻辑设计，规划直行、转弯、碰撞场景方案以及参数对照实验方案。
第四章为系统实现与实验结果分析。基于PhysX5完成仿真代码开发，完成三组交互场景测试，开展多组参数对照实验，记录仿真帧率、形变效果，对比分析PBD与FEM方案的优缺点。
第五章为总结与展望。归纳本文完成的工作，总结实验得出的结论，分析当前仿真系统存在的缺陷，对后续柔性车辆仿真优化方向进行展望。

